\documentclass[10pt,a4paper]{article}

\usepackage[margin=2.0cm]{geometry}
\usepackage[T1]{fontenc}
\usepackage{lmodern}
\usepackage{microtype}
\usepackage{xcolor}
\usepackage{booktabs}
\usepackage{enumitem}
\usepackage{titlesec}
\usepackage{listings}
\usepackage{amsmath}
\usepackage[hidelinks]{hyperref}

\definecolor{rule}{gray}{0.45}
\definecolor{shade}{gray}{0.94}

\titleformat{\section}{\large\bfseries}{\thesection}{0.6em}{}
\titleformat{\subsection}{\normalsize\bfseries}{\thesubsection}{0.6em}{}
\titlespacing*{\section}{0pt}{1.1ex plus .2ex}{0.6ex}

\lstdefinelanguage{Lean4}{
  morekeywords={def,theorem,structure,class,inductive,instance,namespace,end,
    import,open,where,deriving,axiom,example,by,intro,cases,rw,simp,exact,
    unfold,omega,decide,fun,let,in,if,then,else,do,match,with,return,pure,
    private,partial,Type,Prop,Nat,Bool,UInt8,Array,List,Option},
  sensitive=true,
  morecomment=[l]{--},
  morecomment=[s]{/-}{-/},
  morestring=[b]",
}
\lstset{
  language=Lean4,
  basicstyle=\ttfamily\footnotesize,
  columns=fullflexible,
  keepspaces=true,
  breaklines=true,
  frame=single,
  rulecolor=\color{rule},
  framesep=3pt,
  xleftmargin=4pt,
  xrightmargin=4pt,
  aboveskip=0.6em,
  belowskip=0.3em,
  commentstyle=\itshape\color{black!65},
  keywordstyle=\bfseries,
  showstringspaces=false,
  extendedchars=true,
  literate=
    {×}{{$\times$}}1 {→}{{$\to$}}1 {≤}{{$\le$}}1 {∀}{{$\forall$}}1
    {¬}{{$\neg$}}1 {∧}{{$\wedge$}}1 {∨}{{$\vee$}}1 {⟨}{{$\langle$}}1
    {⟩}{{$\rangle$}}1 {α}{{$\alpha$}}1 {β}{{$\beta$}}1 {ρ}{{$\rho$}}1
    {Γ}{{$\Gamma$}}1 {Σ}{{$\Sigma$}}1 {∈}{{$\in$}}1 {≃}{{$\simeq$}}1
    {↦}{{$\mapsto$}}1 {←}{{$\leftarrow$}}1 {·}{{$\cdot$}}1 {∃}{{$\exists$}}1 {∉}{{$\notin$}}1
}
\lstdefinestyle{out}{
  language={},
  basicstyle=\ttfamily\footnotesize,
  frame=none,
  backgroundcolor=\color{shade},
  xleftmargin=8pt,
  aboveskip=0.2em,
  belowskip=0.7em,
}

\newcommand{\file}[1]{\texttt{\small #1}}

\title{\vspace{-1.2cm}\textbf{A capability-ledger subsystem in Lean}\\[2pt]
\large Implementation report, with examples\\[4pt]
\normalsize separate from, and not a claim of, the ITP manuscript}
\author{Gavin Mendel-Gleason}
\date{15 November 2026}

\begin{document}
\maketitle
\vspace{-0.8cm}

\begin{center}
\fbox{\parbox{0.96\textwidth}{\small
\textbf{Scope.} This report describes the \texttt{Cap} subsystem: exact shares,
bounded footprints, a closed value-code universe, an owner/context ledger with
guards, lifetime and loan bookkeeping, a reference byte heap and a pure mutable
backend related to it, a surface checker and certificate command, and a union-find
example under the write discipline. Every Lean fragment is quoted from the
repository and every output was observed at the revision reported. The subsystem
is not a refinement of Lean's compiler or runtime, and the open research stages are
listed in Section~\ref{sec:open}.
}}
\end{center}

\section{What the subsystem is}

Lean has no linear or affine typing, so a ``unique'' wrapper is not unique: a value
can be bound twice. The subsystem therefore separates three things. Handles are
inert names; authority is a position in the current checked context; and the
context is never a public value. Every access is justified against a capability
ledger, and the guarded operations take the ledger's predicates as proofs.

\begin{lstlisting}[caption={Duplication is legal in the host, which is why authority cannot live in a wrapper.}]
structure Owned (α : Type) where
  private mk ::
  val : α

def dup (a : Owned α) : Owned α × Owned α := (a, a)
\end{lstlisting}

\section{Layout}

\begin{center}
\begin{tabular}{@{}ll@{}}
\toprule
Module & Content \\
\midrule
\file{Cap/Id.lean} & allocation, capability, loan, lifetime identities \\
\file{Cap/Share.lean} & shares, splits, the one-transfer theorem \\
\file{Cap/Footprint.lean} & half-open intervals, views, split facts \\
\file{Cap/Code.lean} & codes for the closed value universe \\
\file{Cap/Admission.lean} & bijection certificates and the encoding gate \\
\file{Cap/Ledger.lean} & decomposition tree, budget, owners, splits, joins, guards \\
\file{Cap/Context.lean} & multi-owner context, dispatch, freshness, disposal \\
\file{Cap/Loans.lean} & lifetime and loan bookkeeping \\
\file{Cap/Borrow.lean} & the integrated delegation invariant \\
\file{Cap/Discipline.lean} & world and delegation-chain correspondence \\
\file{Cap/Delegation.lean} & nested delegation chain, rooted path, checked begin/end \\
\file{Cap/Reconstruction.lean} & borrowed subdivision and reconstruction obligations \\
\file{Cap/Integrated.lean} & integrated multi-allocation sealed execution \\
\file{Cap/UFReuse.lean} & reusable two-find compression on the functional model \\
\file{Cap/UFIntegrated.lean} & two full-path finds through the integrated sealed boundary \\
\file{llvm-bridge/} & compiled-C store kernel and its composition with an authorized write \\
\file{Cap/FullInv.lean} & the combined invariant and its preservation \\
\file{Cap/Bytes.lean} & reference byte heap \\
\file{Cap/Backend.lean} & pure mutable backend and its refinement \\
\file{Cap/BackendRun.lean} & the 64-byte run on the backend \\
\file{Cap/Exec.lean} & instruction interpreter over the backend \\
\file{Cap/World.lean} & owners plus mutable backends \\
\file{Cap/Core.lean} & deep command language and executable checker \\
\file{Cap/Surface.lean} & surface checker and certificate command \\
\file{Cap/Certify.lean} & a kernel-checked closed-program certificate \\
\file{Cap/UnionFind.lean} & union-find under the write discipline \\
\file{Cap/Unified.lean} & value-carrying writes, checker as interpreter \\
\file{Cap/Disposal.lean} & the legal-disposal predicate and checked free \\
\bottomrule
\end{tabular}
\end{center}

\section{Changes since the reviewed revision}

\begin{center}
\begin{tabular}{@{}lp{0.60\textwidth}@{}}
\toprule
Review item & Change in this revision \\
\midrule
R1 & \file{Cap/Unified.lean}: a value-carrying command language --- write, both splits, join, free, alloc, borrow/return and no-op --- with the checker as interpreter and the world relation preserved; the 64-byte and lifecycle runs go through it. Observations are taken from the post-world; a snapshot command is not a command of this language. \\
R2 & Access-capability ids separated from internal split-node provenance; \texttt{Owner.FreshIds} and a freshness branch in the checker refusing duplicate and reused-parent identifiers (\texttt{msgFresh}); regressions in \texttt{bench/CapFreshness.lean}. The full invariant is still open. \\
R3 & Delegation by retagging is the borrowed-authority mechanism: a borrow retags the active leaf to a distinct borrower and records a loan with its chain base, a return retags the borrower back and refuses while a descendant loan is live, and \texttt{World.BorrowOK} carries distinct names, issuance and liveness, the recorded base, global lender detachment, active-terminal-or-delegated authority, adjacent-edge agreement and issued-order monotonicity. \texttt{World.checkWorld\_fullInv} preserves it; \texttt{World.delegationChains\_ok} proves the world/chain correspondence under the separate \texttt{World.LoansChainOrderOK} hypothesis. \\
R4 & \file{Cap/Disposal.lean}: one \texttt{Legal} predicate (entitlement, coverage, reconstructed root, empty loans) and a checked \texttt{free?}; not yet a command of the unified language. \\
R5 & Not done: the complete split/borrow/fill/read/rejoin/restore/snapshot/free trace. \\
G7 & \file{Cap/UFReuse.lean}: two full-path finds on the functional model, \texttt{[0,0,1,2,2]} from node 3 to \texttt{[0,0,0,0,2]} then from node 4 to \texttt{[0,0,0,0,0]}, with the first output's ranked-forest witness admitting the second operation as a certificate rather than a restated array. \file{Cap/UFIntegrated.lean} now runs the same two calls through the integrated sealed boundary on the first call's returned state, writing through the whole-share capability and re-exporting the integrated invariant; the world-based \texttt{UFRunner} is still not threaded. \\
Accounting & \texttt{Reconstruction.State.Inv} now includes the ownership account \texttt{State.Accounts}, which relates the forest to \texttt{issued}, \texttt{retired} and \texttt{restores}: current authority-node identities, issuance coverage, retirement, dormant lenders, restoration realization, resource agreement and nested-obligation agreement. The integrated invariant requires the session-persistent allocation and capability histories, and \texttt{free} leaves both intact; \texttt{State.alloc} refuses an allocation identifier or an initial capability name that was ever issued, so a freed allocation identifier cannot be reissued and an old capability cannot authorize a replacement. \\
Kernel & \file{llvm-bridge/}: clang's emitted body for one-byte \texttt{store8} is proved to have exactly the reference store's bounded outcome, and the composed theorem relates one accepted integrated write to that run at a caller-supplied address window, preserving the post-state's representation, the frame and the integrated invariant. The address assignment, IR projection, ABI and native execution remain stated assumptions. \\
Reconstruction & \file{Cap/Reconstruction.lean} records checked spatial and fractional subdivision rules, a reconstruction operation and restoration obligations, each preserving the reconstruction invariant including the ownership account; \file{Cap/Integrated.lean} lifts it to a multi-allocation state and \file{Cap/Sealed.lean} re-exports the integrated sealed boundary, whose invariant includes each component's account and storage agreement. The world \texttt{UCmd} split/join guard remains allocation-wide. \\
Validation & \texttt{scripts/check\_cap\_controls.sh} requires each negative fixture to fail with its intended diagnostic; \texttt{check\_all} runs it. \\
Examples & Union-find and a linear-probing hash set under the guard API, with their limits stated. \\
\bottomrule
\end{tabular}
\end{center}

\section{Exact shares and footprints}

A context fixes a denominator $D$; a share is a positive numerator $q \le D$. A
write needs the whole unit, a read only a positive share.

\begin{lstlisting}[caption={\file{Cap/Share.lean}: the one-transfer theorem.}]
def ShareWf (D q : Nat) : Prop := 0 < q ∧ q ≤ D
def Splits (q q1 q2 : Nat) : Prop := q1 + q2 = q
def Writable (D q : Nat) : Prop := q = D

theorem no_two_writes {D q q1 q2 : Nat} (h : Splits q q1 q2) (h1 : 0 < q1)
    (hw1 : Writable D q1) (hw2 : Writable D q2) (hq : q ≤ D) : False
\end{lstlisting}

Footprints are half-open intervals with an ordering proof; a view names an
allocation and a footprint. Splitting at a point yields disjoint children that
cover the parent, and an allocation bound passes through:

\begin{lstlisting}[caption={\file{Cap/Footprint.lean}, selected facts.}]
theorem splitAt_disjoint ... : (splitAt i k h1 h2).1.Disjoint (splitAt i k h1 h2).2
theorem splitAt_adjacent ... : (splitAt i k h1 h2).1.Adjacent (splitAt i k h1 h2).2
theorem splitAt_cover ... :
    (splitAt i k h1 h2).1.lo = i.lo ∧ (splitAt i k h1 h2).2.hi = i.hi
theorem splitAt_within ... : Within len (...) ∧ Within len (...)
\end{lstlisting}

\section{Admission: encoding is not confinement}

A type is admitted to the \emph{value language} by a demonstrated bijection to a
closed code. The gate reflects the named demonstration, requires the code to be
closed, and audits the axioms.

\begin{lstlisting}[caption={\file{Cap/Admission.lean}.}]
structure Bij (α β : Type) where
  toFun : α → β
  invFun : β → α
  left_inv : ∀ a, invFun (toFun a) = a
  right_inv : ∀ b, toFun (invFun b) = b

structure SafeValue (T : Type) where
  code : Code
  bij : Bij T (El code)
\end{lstlisting}

\begin{lstlisting}[style=out]
'Pair' is encoded by code Code.nat.prod Code.nat
the demonstration 'badBij' depends on non-standard axioms
  [badFrom, badLeft, badRight, badTo]
the demonstration 'Other.bij' does not have type Bij Pair (El c)
\end{lstlisting}

A bijection is not a confinement proof. \texttt{Unit → Nat} is bijective with
\texttt{El .nat} and a particular thunk can capture a buffer:

\begin{lstlisting}[caption={\file{bench/NatThunkAdmission.lean}.}]
abbrev NatThunk := Unit → Nat

def thunkBij : Bij NatThunk (El .nat) where
  toFun f := f ()
  invFun n := fun _ => n
  left_inv f := by funext u; cases u; rfl
  right_inv n := rfl

def retainBuffer (bytes : Array UInt8) : NatThunk := fun _ => bytes.size
\end{lstlisting}

The gate accepts the encoding and the storage type excludes it: \texttt{SafeType}
holds a code and an initialiser producing \texttt{El code}, never an arbitrary
\texttt{T}.

\section{The ledger}

Capabilities are leaves of a decomposition tree. Internal nodes are fractional and
spatial splits; the budget at a byte is the frontier sum. Splits and sibling joins
preserve the budget:

\begin{lstlisting}[caption={\file{Cap/Ledger.lean}.}]
theorem frac_split_budget ... :
    (Tree.frac c1.id (.leaf c1) (.leaf c2)).budget x = (Tree.leaf c).budget x

theorem spatial_split_budget ... :
    (Tree.spatial c1.id (.leaf c1) (.leaf c2)).budget x = (Tree.leaf c).budget x

theorem Tree.budget_replace ... :
    (t.replace id f).budget x = t.budget x
\end{lstlisting}

At the owner level, \texttt{Owner.Covered} is the byte-wise full-budget invariant;
splits preserve it and allocation establishes it. The write guard is a
proposition:

\begin{lstlisting}[caption={\file{Cap/Ledger.lean}: coverage and guards.}]
def Owner.Covered (o : Owner) (D : Nat) : Prop :=
  ∀ x, x < o.len → o.root.budget x = D

def Owner.Writable (o : Owner) (D : Nat) (id : CapId) (i : Nat) : Prop :=
  ∃ c, o.root.find? id = some c ∧ c.state = .active ∧ c.q = D ∧
    c.ceiling = .readWrite ∧ c.lo ≤ i ∧ i < c.hi

theorem Owner.splitFraction_covered ... : o'.Covered D
theorem Owner.splitAt_covered ... : o'.Covered D
\end{lstlisting}

Capability identifiers are collected from the tree's leaves only; internal split
nodes record provenance and are excluded. A split must use \texttt{Owner.FreshIds}:
two distinct child identifiers absent from the owner, which includes the parent.
The deep checker enforces this with \texttt{msgFresh}, so duplicate children and a
reused parent identifier are refused:

\begin{lstlisting}[caption={\file{Cap/Core.lean}: freshness.}]
theorem splitAt_checked_fresh ... : o.FreshIds id1 id2
theorem splitFraction_checked_fresh ... : o.FreshIds id1 id2
\end{lstlisting}

The join is sibling-only, refuses aliases, and preserves coverage:

\begin{lstlisting}[caption={\file{Cap/Join.lean}.}]
theorem Owner.joinFraction_covered ... : o'.Covered D
theorem Owner.joinFraction_alias_none ... :
    -- a join whose two children carry the same capability id is refused
\end{lstlisting}

\section{Context, lifecycle and disposal}

The context is a list of owners with distinct allocations. Allocation appends a
fresh fully covered owner, disposal removes one through its recorded disposal
capability, and splits dispatch to the owner:

\begin{lstlisting}[caption={\file{Cap/Context.lean}.}]
def Ctx.alloc (C : Ctx) (a : AllocId) (n : Nat) (pos : 0 < C.D)
    (rootId disposal : CapId) (hfresh : a ∉ C.owners.map (·.alloc)) : Ctx

def Ctx.freeWith? (C : Ctx) (a : AllocId) (d : CapId) : Option Ctx
def Ctx.splitAtOwner? ... : Option Ctx
def Ctx.splitFractionOwner? ... : Option Ctx

theorem Ctx.alloc_realizes ... : (C.alloc ...).Realizes
theorem Ctx.freeWith?_realizes ... : C'.Realizes
theorem Ctx.splitAtOwner_realizes ... : C'.Realizes
\end{lstlisting}

Lifetimes and loans are pure bookkeeping; the retag accounting model turns them
into delegation. A borrow moves the whole authority to a distinct borrower, and
the world invariant records the chain:

\begin{lstlisting}[caption={\file{Cap/Borrow.lean}, \file{Cap/Discipline.lean}, selected facts.}]
def World.BorrowOK (w : World) : Prop :=
  LoansDistinctOK w ∧ LoansNonAliasOK w ∧ LoansIssuedOK w ∧ LifetimesBoundedOK w
  ∧ LoansAllocOK w ∧ LoansDetachedOK w ∧ LoansBaseOK w ∧ LoansTerminalOK w
  ∧ LoansAgreeOK w ∧ LoansOrderOK w ∧ w.BorrowAcyclic

theorem World.delegationChains_ok (h : w.BorrowOK) (hord : w.LoansChainOrderOK) :
    -- every chain is valid; the chains cover every loan; each chain's
    -- projection is its (allocation, base) subsequence; the flattened
    -- projection is a permutation of the world loan list
\end{lstlisting}

\texttt{LoansChainOrderOK} is a separate hypothesis: \texttt{BorrowOK} does not
determine the order of the world loan list.

\section{The 64-byte acceptance run}

The concrete unified lifecycle allocates a 4-byte owner, writes, splits, writes,
splits fractionally, joins, writes and disposes; the pre-existing 64-byte backend
is untouched:

\begin{lstlisting}[caption={\file{Cap/Unified.lean}.}]
theorem uLifecycle_snapshot :
    uLifecycleSnapshot = some (some (List.replicate 64 (0 : UInt8))) := by decide
\end{lstlisting}

The reference heap proves write visibility, length preservation, non-interference
at other indices, snapshot independence, and that two readers of the right half each
see 64:

\begin{lstlisting}[caption={\file{Cap/Bytes.lean}, selected facts.}]
theorem write_get ... : read? h' a i = some v
theorem write_preserves_length ... : ...
theorem write_ne ... : read? h' a j = read? h a j
theorem right_readers_both_64 : rightReaders = some (64, 64)
theorem snapshot_independent : snapshotRun = some ([7,7,7,7], [7,7,7,7], [9,7,7,7])
\end{lstlisting}

\section{Backend, interpreter and world}

A pure mutable backend refines the reference heap; the 64-byte run on it agrees;
an instruction interpreter agrees on the example; the world adds owners with
backends and preserves the relation under guarded write and disposal.

\begin{lstlisting}[caption={\file{Cap/Backend.lean}, \file{Cap/BackendRun.lean}, \file{Cap/Exec.lean}.}]
theorem Backend.read_refines ... : h.read? a i = some v
theorem Backend.write_refines ... : b'.Rel h' a
theorem Backend.write_then_read ... : b'.read i = some v

theorem exampleBackend_snapshot :
    exampleBackendSnapshot = some (List.replicate 32 (1:UInt8) ++ List.replicate 32 (2:UInt8))

theorem exec_example :
    exec (Backend.empty 64) execExampleProgram = some (_, List.replicate 32 1 ++ List.replicate 32 2)
\end{lstlisting}

\begin{lstlisting}[caption={\file{Cap/World.lean}.}]
theorem World.write_rel ... : w'.Rel
theorem World.freeWith_rel ... : w'.Rel
\end{lstlisting}

\section{The stronger owner and context invariant}

\file{Cap/Invariant.lean} strengthens the ledger invariant with the facts the
runner handoff asks for at this layer: every leaf of an owner's tree lies inside
the owner's allocation and extent, coverage and the recorded disposal capability
hold, and a context has realization, per-owner well-formedness and distinct
capability identifiers.

\begin{lstlisting}[caption={\file{Cap/Invariant.lean}.}]
def Owner.LeavesWithin (o : Owner) : Prop :=
  forall id c, o.root.find? id = some c -> c.alloc = o.alloc ∧ c.hi <= o.len

def Owner.WellFormed (o : Owner) (D : Nat) : Prop :=
  o.Covered D ∧ o.LeavesWithin ∧ o.DisposalPresent

def Ctx.WellFormed (C : Ctx) : Prop :=
  C.Realizes ∧ (forall o ∈ C.owners, o.WellFormed C.D) ∧ C.ids.Nodup

theorem Owner.init_wellFormed ...
theorem Ctx.alloc_wellFormed ...
theorem Ctx.splitAtOwner_wellFormed ...
theorem Ctx.splitFractionOwner_wellFormed ...
theorem Ctx.free_wellFormed ...
\end{lstlisting}

Establishment and preservation are proved with explicit hypotheses: owner id
distinctness and freshness for the owner splits, context-wide child freshness for
the dispatched splits, and allocation and id freshness for \texttt{alloc}. The tree's
first-occurrence replacement is characterised (\texttt{find?\_replace\_eq},
\texttt{ids\_replace\_perm}, \texttt{leaves\_replace\_nodup}) and the identifier list is shown to be
preserved up to permutation. The concrete 64-byte owner satisfies `WellFormed 64`.

Not yet in this invariant: loans and lifetime association, storage/heap agreement,
one disposal entitlement per allocation as a global condition, and ceiling
inheritance. \texttt{DisposalPresent} is deliberately a disjunction (a leaf, or absence
from the tree's ids) so that initialisation is unconditional; the strict reading is
recorded separately.

\section{The world-level invariant}

\file{Cap/WorldInvariant.lean} bundles the three components needed above the
owner/context layer: the context invariant, owner uniqueness per allocation, and
storage agreement.

\begin{lstlisting}[caption={\file{Cap/WorldInvariant.lean}.}]
def Ctx.Functional (C : Ctx) : Prop :=
  forall a o o', C.findOwner? a = some o -> o' ∈ C.owners -> o'.alloc = a -> o' = o

def World.StorageMatches (w : World) : Prop :=
  forall a o, w.ctx.findOwner? a = some o ->
    exists b, w.heaps a = some b ∧ b.buf.length = o.len

def World.WellFormed (w : World) : Prop :=
  Ctx.WellFormed w.ctx ∧ Ctx.Functional w.ctx ∧ w.StorageMatches

theorem World.rel_of_wellFormed : w.WellFormed -> w.Rel
theorem World.write?_wellFormed ... : w'.WellFormed
theorem World.freeWith?_wellFormed ... : w'.WellFormed
\end{lstlisting}

Functionality is derived from distinct allocation identifiers
(\texttt{functional\_of\_nodup}); update and free preserve it; write and disposal
preserve the bundle with no additional hypotheses. Loans, global disposal
uniqueness and the split commands are not yet part of this layer.

\section{One legal disposal rule}

\file{Cap/Disposal.lean} gives the single predicate under which an allocation may
be freed. A `DisposalState` pairs the world with the loan context; `Legal` requires
the found owner, its coverage, the recorded disposal entitlement, an empty loan
list, and a reconstructed root (a single active whole-share leaf over the
allocation). `free?` is the checked route through the world's disposal check.

\begin{lstlisting}[caption={\file{Cap/Disposal.lean}.}]
def DisposalState.Legal (s : DisposalState) (a : AllocId) (d : CapId) : Prop :=
  exists o, s.world.ctx.findOwner? a = some o ∧ o.Covered s.world.ctx.D ∧
    o.disposal = d ∧ s.loans.loans = [] ∧ o.RootReconstructed s.world.ctx.D

def DisposalState.free? ... : Option DisposalState

theorem DisposalState.free?_none_of_not_legal ... : s.free? a d = none
theorem DisposalState.free?_legal ... : exists s', s.free? a d = some s'
theorem DisposalState.free?_removes_owner ... : s'.world.ctx.findOwner? a = none
theorem DisposalState.free?_none_after_free ... : s'.free? a d = none
\end{lstlisting}

Refusals for a wrong entitlement and for an outstanding loan are proved, and a
legal disposal succeeds and leaves no owner for the allocation. The predicate is
not yet wired into `Cap.Unified`; names are invalidated by owner removal rather
than by arena generations.

\section{Unified certified execution (first slice)}

The review's principal finding was that certification and execution were separate.
The first slice unifies them for byte writes: a command carries the allocation, the
capability, the index and the value; the checker is the interpreter over the world;
and acceptance names the guard while preserving the world relation.

\begin{lstlisting}[caption={\file{Cap/Unified.lean}.}]
inductive UCmd where
  | write (a : AllocId) (id : CapId) (i : Nat) (v : UInt8)
  | nop

def World.writeChecked (w : World) (a : AllocId) (id : CapId) (i : Nat) (v : UInt8) :
    Except String World :=
  match w.ctx.findOwner? a with
  | none => .error "write: no owner for the allocation"
  | some o =>
    match w.heaps a with
    | none => .error "write: no backend for the allocation"
    | some b =>
      if _hw : o.writableB w.ctx.D id i then
        match b.write i v with
        | none => .error "write: index out of range"
        | some b' =>
          .ok { w with heaps := fun x => if x = a then some b' else w.heaps x }
      else .error "write: capability is not writable at this index"

def checkWorld : World -> List UCmd -> Except String World
  -- checker and interpreter: an accepted program returns the world it executed

theorem World.writeChecked_sound ... : -- names the owner, its guard, the backend and the replacement
theorem World.writeChecked_rel ... : w'.Rel
theorem checkWorld_rel ... : w'.Rel
\end{lstlisting}

\begin{lstlisting}[caption={The 64-byte run through the certified route.}]
def uProgram : List UCmd := fillCmds ⟨1⟩ 0 32 1 ++ fillCmds ⟨1⟩ 32 32 2

theorem uRun_snapshot :
    uRun = some (List.replicate 32 (1 : UInt8) ++ List.replicate 32 (2 : UInt8)) := by
  decide

theorem uProgram_rel : exists w, checkWorld uWorld uProgram = .ok w ∧ w.Rel
\end{lstlisting}

Scope: the command language covers writes, both splits, join, allocation,
borrow/return, disposal and no-op. The surface still targets the owner-layer
checker.

\section{World-indexed computations}

\file{Cap/Indexed.lean} gives the computation type the design calls for: a
program is indexed by its pre-world and by a function from its result to its
post-world, and composition forces the continuation to start in the world the
first computation produced.

\begin{lstlisting}[caption={\file{Cap/Indexed.lean}.}]
structure Ix (pre : World) (α : Type) (post : α -> World) where
  run : (w : World) -> w = pre -> Option α

def Ix.bind (m : Ix pre α post) (k : (a : α) -> Ix (post a) β post') :
    Ix pre β post'
  -- the continuation's pre-world is the first computation's post-world

def writeIx ... : Ix pre Unit (fun _ => pre.writePure a b')
  -- carries the ledger guard and the backend write; the post-index is the world
  -- the guarded write produces

def prog : Ix uWorld Unit (fun _ => post2)
theorem prog_run : prog.run uWorld rfl = some ()
theorem prog_sound :
    checkWorld uWorld [UCmd.write ..., UCmd.write ...] = .ok post2
\end{lstlisting}

The two-write program is checked at elaboration: \texttt{prog\_run} computes its run and
\texttt{prog\_sound} shows the executable route reaches the same post-world. The indices
are ordinary `World` values; there is no elaborator inferring them for symbolic
programs, and runtime mutation is still performed by the executable route.

\section{Sealed program certificates}

The runner handoff's sharpest point is that a type with a bare `run` field does
not establish a post-relation: a client can inhabit it with an unrelated function.
\file{Cap/Runner.lean} gives the certificate-carrying form: syntax packaged with
the checker's derivation.

The closed form starts from the canonical empty world (no owners, no backends,
well formed by \texttt{emptyWorld\_wf}); a \texttt{ClosedProgram} packages
commands accepted from it, and the only observation is the snapshot of the final
well-formed world. \texttt{closed64} is the allocate/fill example, with
\texttt{closed64\_snapshot} proving the 32 ones then 32 twos.

\begin{lstlisting}[caption={\file{Cap/Runner.lean}.}]
structure Prog (pre : World) where
  post : World
  cmds : List UCmd
  ok : checkWorld pre cmds = .ok post

structure ClosedProgram where
  cmds : List UCmd
  accepted : (checkWorld World.emptyWorld cmds).isOk = true

def Prog.snapshot (p : Prog pre) (a : AllocId) : Option (List UInt8) :=
  (p.post.heaps a).map (fun b => b.buf)

def p64 : Prog uWorld := ⟨post64, uProgram, post64_ok⟩
theorem p64_snapshot : p64.snapshot ⟨0⟩ = some (32 ones ++ 32 twos)
\end{lstlisting}

A `Prog` cannot exist without the checker's result, and the only observation is an
immutable snapshot from the post-world. Certificates compose: \texttt{checkWorld\_append}
runs an appended command list from the first part's result, and \texttt{Prog.seq} makes
the second program start in the first one's post-world. \texttt{Cap/Indexed.lean}'s `Ix` is a
specification type by comparison: its `run` field alone does not establish a
post-relation, which is now stated in that module. \file{Cap/Sealed.lean} re-exports the integrated sealed boundary from
\file{Cap/Integrated.lean}: \texttt{Sealed} is \texttt{Integrated.Sealed}, with a
private constructor, a private representation field and a private invariant
proof. The canonical state is one whole fresh allocation
(\texttt{Sealed.initial}); \texttt{Sealed.run} executes the checked integrated
\texttt{Cmd} list and carries \texttt{State.run\_inv} through
(\texttt{Sealed.run\_invariant}). Negative fixtures check that the constructor,
the representation field and the retired world entry point are unavailable
(\texttt{bench/CapSealReject.lean}, \texttt{bench/CapSealPeekReject.lean}). The
observations are a per-allocation snapshot, byte read, fragment description,
retired access names and outstanding restoration obligations.
\texttt{boundaryClaims} lists the executable-edge claims with their status:
integrated command execution, invariant preservation and the observation
interface are proved, while the C runtime's store trace, physical reference
counts and physical storage independence are runtime contracts and concurrency is
excluded. One clang-emitted body's trace is proved in \file{llvm-bridge/} and
composed with an authorized integrated write (Section~\ref{sec:kernel}); that is
not the runtime contract and not a physical ownership refinement.
Physical storage independence is not a theorem of the pure model --- copied
values share their representation --- so the C runner's contract remains open.
The invariant the sealed state
carries is the integrated state invariant: distinct live allocation identifiers,
each component's reconstruction invariant with its ownership account, each
component's storage agreement
(\texttt{Reconstruction.State.Storage}), and the session-persistent allocation
and capability histories. \texttt{World.FullInv} remains a
separate invariant on the world \texttt{UCmd} path, not the sealed state's
carrier.
\texttt{World.BorrowOK} is the integrated delegation invariant: distinct
loan/lender/borrower/lifetime names, issuance and liveness, the recorded chain
base, global lender detachment, active-terminal-or-delegated authority,
adjacent-edge agreement and issued-order acyclicity.

\section{Surface checker and certificates}

\texttt{\#own\_prog} elaborates surface commands to the deep language and runs the
checker at elaboration time; \texttt{\#own\_cert} declares the kernel-checked
theorem for a closed program.

\begin{lstlisting}[caption={\file{Cap/Surface.lean}, surface program and certificate.}]
#own_prog surfaceOwner 64 {
  splitAt ⟨0⟩ ⟨2⟩ ⟨3⟩ 32
  write ⟨2⟩ 5
  write ⟨3⟩ 40
  splitFraction ⟨2⟩ ⟨4⟩ ⟨5⟩ 32 32
  nop
}

#own_cert surfaceCert : surfaceOwner 64 {
  splitAt ⟨0⟩ ⟨2⟩ ⟨3⟩ 32
  write ⟨2⟩ 5
  write ⟨3⟩ 40
}
\end{lstlisting}

\begin{lstlisting}[style=out]
capability program accepted
Cap.surfaceCert :
  Checked surfaceOwner 64
    [Cmd.splitAt { val := 0 } { val := 2 } { val := 3 } 32,
     Cmd.write { val := 2 } 5, Cmd.write { val := 3 } 40]
\end{lstlisting}

The certificate is axiom-free; the deep checker has write and split soundness and
carries coverage through a run:

\begin{lstlisting}[caption={\file{Cap/Core.lean}, \file{Cap/Certify.lean}.}]
theorem write_checked ... : o.Writable D id i
theorem check_covered ... : o'.Covered D
theorem certifiedProgram_checked : Checked certifiedOwner 12 certifiedProgram
\end{lstlisting}

\section{Union-find under the discipline}

The parent and rank arrays are one allocation, covered by a single whole-share
read--write leaf. Every write is \texttt{UF.writeParent}, which takes the guard as a
proof and preserves the store length. In-range guards are supplied; index 4 has
none.

\begin{lstlisting}[caption={\file{Cap/UnionFind.lean}.}]
structure UF where
  owner : Owner
  parent : List Nat
  rank : List Nat
  hlen : parent.length = rank.length
  D : Nat

def UF.writeParent (u : UF) (id : CapId) (i v : Nat)
    (_hg : u.owner.Writable u.D id i) : UF :=
  { owner := u.owner, parent := u.parent.set i v, rank := u.rank,
    hlen := by rw [List.length_set]; exact u.hlen, D := u.D }

theorem ufOwner_writable (i : Nat) (h : i < 4) : ufOwner.Writable 4 ⟨1⟩ i
theorem ufOwner_not_writable_4 : ¬ ufOwner.Writable 4 ⟨1⟩ 4
\end{lstlisting}

The concrete run attaches 1 under 0, 3 under 2, and 2 under 0:

\begin{lstlisting}[caption={The run and its results.}]
def uf0 : UF :=
  { owner := ufOwner, parent := [0, 1, 2, 3], rank := [0, 0, 0, 0], hlen := rfl, D := 4 }

def uf1 : UF := uf0.writeParent ⟨1⟩ 1 0 (ufOwner_writable 1 (by decide))
def uf2 : UF := uf1.writeParent ⟨1⟩ 3 2 (ufOwner_writable 3 (by decide))
def uf3 : UF := uf2.writeParent ⟨1⟩ 2 0 (ufOwner_writable 2 (by decide))

theorem uf3_parent : uf3.parent = [0, 0, 0, 2] := by decide

theorem uf3_find_all :
    uf3.findRoot 0 = some 0 ∧ uf3.findRoot 1 = some 0 ∧
      uf3.findRoot 2 = some 0 ∧ uf3.findRoot 3 = some 0 := by decide
\end{lstlisting}

What the example establishes: every write in the run is authorised by the ledger
guard, out-of-range writes have no guard, the parent list is as stated, all four
lookups reach 0, and the arrays keep their lengths. What it does not establish: a
union-find correctness theorem (partition refinement, path-compression semantics,
complexity), or rank balancing --- the example attaches by explicit choice.

\subsection*{What the examples enforce, and what they do not}

The hash set in \file{Cap/HashSet.lean} has the same status. Both examples require
a \texttt{Writable} proof for every write, supplied from a whole-share owner, so
the write is proof-carrying at the owner layer. But the proof is not yet
load-bearing: in both write functions the guard parameter is unused; neither
structure records that the store's length and allocation agree with the owner
(\texttt{parent.length = owner.len}, \texttt{slots.length = owner.len}); and
neither algorithm runs through \texttt{World.writeChecked}/\texttt{checkWorld}, so
\texttt{World.Rel} is not maintained for them. A caller could pair a writable owner
with a store of another length or allocation, and \texttt{List.set} would silently
do nothing out of range rather than being refused. Closing this needs the store
relation, the bounded-leaves fact, a \texttt{write\_in\_bounds} theorem that uses
the guard, and execution through the unified route.

\section{A second example: a linear-probing hash set}

\file{Cap/HashSet.lean} puts a hash set under the write discipline: one eight-slot
whole-share allocation, insertion probing from the hash position and writing
through the owner's \texttt{Writable} guard, a read-only lookup, and
\texttt{insert\_length} preserving the table length.

\begin{lstlisting}[caption={\file{Cap/HashSet.lean}.}]
def Table.write (t : Table) (i : Nat) (v : Option Nat)
    (_hg : t.owner.Writable t.D ⟨1⟩ i) : Table

def Table.insert (t : Table) (x : Nat) (howner : t.owner = hsOwner) (hD : t.D = 8)
    : Nat -> Nat -> Table

theorem Table.insert_length ... :
    (t.insert x howner hD fuel start).slots.length = t.slots.length
\end{lstlisting}

The concrete colliding insertions 3, 11 and 19 fill slots 3, 4 and 5; membership
of 11 is true, of 4 false, and the length stays 8. No resizing, deletion, hashing
theorem or complexity claim is made.

\section{Union-find partition semantics}

\file{Cap/UFPartition.lean} is the algorithmic relation, separate from ownership:
\texttt{Reachable} follows parent steps, \texttt{Same} holds when two nodes reach
a common root and is proved reflexive, symmetric and transitive, and the simple
union links one node directly under another.

\begin{lstlisting}[caption={\file{Cap/UFPartition.lean}.}]
inductive Reachable (p : Arr) : Nat -> Nat -> Prop
  | one {a b} : step p a = b -> Reachable p a b
  | more {a b c} : step p a = b -> Reachable p b c -> Reachable p a c

def Same (p : Arr) (a b : Nat) : Prop :=
  exists r, Reachable p a r ∧ Reachable p b r

theorem same_refl, same_symm, same_trans
theorem union_same ... : Same (union p a b) a b
theorem uf3_eq : uf3 = [0, 0, 0, 2]
theorem uf3_same_all : -- every pair below 4 is Same
\end{lstlisting}

This is the relation, not the algorithm library: the union is the direct
\texttt{set b a}, with no root finding, rank, path compression or ownership
integration, and no complexity claim. It supplies the partition facts an
ownership-integrated union-find would need to preserve.

\section{Root finding and compression}

\file{Cap/UFCompress.lean} adds bounded root finding and one-step compression to
the functional model.

\begin{lstlisting}[caption={\file{Cap/UFCompress.lean}.}]
def findRootF (p : Arr) : Nat -> Nat -> Nat
def StopsAt (p : Arr) (fuel start : Nat) : Prop
theorem findRootF_root ... : IsRoot p (findRootF p fuel start)
theorem findRootF_reachable ... : Reach p start (findRootF p fuel start)
def compress (p : Arr) (x r : Nat) : Arr := p.set x r
theorem compress_same_sound ... : -- compression from a reachable root creates no new Same pairs
\end{lstlisting}

This is one-step compression on the functional model: not full path compression,
not the ownership layer, and no termination or complexity claim. It supplies the
algorithmic facts an ownership-integrated union-find would preserve.

\section{World-level invariant and the split command}

\file{Cap/WorldInvariant.lean} bundles the world invariant: the context invariant
of \file{Cap/Invariant.lean}, one owner per allocation (functionality), and the
agreement that the found owner has a backend of its length. Writes and disposal
preserve it. The spatial split command is at the world level:
\texttt{World.splitAtChecked} requires the owner and both children to be fresh
for the context and preserves the invariant (\texttt{World.splitAtChecked\_wf});
the heap agreement through a same-length owner replacement is
\texttt{World.update\_storage}. The unified checker now carries the invariant:
\texttt{checkWorld\_wf} shows an accepted execution preserves it through writes
and splits, and \texttt{checkWorld\_rel} derives the older \texttt{World.Rel}
from it. \file{Cap/Trace.lean} lifts the per-write report to traces: in any
accepted command list, a write at any position held its guard at the world
reached just before it (\texttt{checkWorld\_write\_guard}), via the single-step
decomposition \texttt{checkWorld\_cons\_ok}. The language also carries the sibling-only fractional split and join and disposal
and allocation commands, each preserving the invariant; the join's context-side
preservation is \texttt{Ctx.joinFractionOwner\_wellFormed}, using an exact
permutation of the capability ids (the two child ids removed, the fresh id added).
The concrete lifecycle program runs alloc, write, spatial split, write, fractional
split, join, write and free through the unified route and leaves the pre-existing
backend at its initial zeros (\texttt{uLifecycle\_snapshot}).
Splits and fractional joins are refused while any loan is outstanding on the
allocation, so the world route does not consult the reconstruction forest
(Section~\ref{sec:reconstruction}); \texttt{World.splitAtChecked\_refuses\_with\_loan}
records that allocation-wide guard.

Borrowing is the retag discipline: \texttt{beginBorrow} retags the lender's
active leaf to a distinct borrower and records the loan, with its chain base, in
the world's loan context, and \texttt{endBorrow} retags the terminal borrower
back to the lender and removes the record. While the loan is outstanding a write
through the lender is refused by the same leaf-name guard
(\texttt{uLoan\_borrower\_write\_refused}, with the accepted write before the
borrow as positive control), and after the return the same write is accepted
(\texttt{uLoanRun\_ok}); \(\texttt{borrowOK\_beginBorrow}\) and
\(\texttt{borrowOK\_endBorrow}\) preserve \texttt{World.BorrowOK}, and
\texttt{World.checkWorld\_fullInv} preserves the combined invariant. A snapshot
command is not yet in this language.

\section{Union-find through the checked world}

\file{Cap/UFRun.lean} runs union-find as a client of the checked execution. Reads
are observations of the pure backends (the discipline constrains writes), and
\texttt{union} issues its writes through \texttt{World.writeChecked}, so
\texttt{union\_wf} preserves the world invariant. The concrete trace attaches 1
under 0, 3 under 2 and 2 under 0; it is accepted (\texttt{ufTrace\_ok}), its
final world is well formed, its backend parent array is exactly the functional
model's \texttt{[0,0,0,2]} (\texttt{ufTrace\_backend}), and every pair of the
four nodes is \texttt{Same} in the functional partition relation
(\texttt{ufTrace\_partition}). Path compression in the client and a read
discipline are not claimed; the compression facts the client would preserve live
in \file{Cap/UFCompress.lean}.

\file{Cap/UFChecked.lean} supplies the read discipline. \texttt{UFRep w a id n}
is the representation predicate of a whole-share read-write parent array whose
stored bytes are all valid node indices. From it, a runtime index \texttt{i < n}
yields the static evidence (\texttt{UFRep.canRead}, \texttt{UFRep.canWrite}), a
represented read is proved to return a value below \texttt{n}
(\texttt{readParent\_lt}) and below 256, and \texttt{findRootC} is the chase
typed by exactly that evidence. \texttt{compressC} is the one-step checked
compression, \texttt{compressC\_ok} says it is never refused and
\texttt{compressC\_wf} preserves the world invariant. Semantically, the backend
is the functional model's array (\texttt{UFArr}); a read is the model's step
(\texttt{readParent\_step}), a completed chase is \texttt{UF.findRootF}
(\texttt{findRootC\_eq\_of\_stops}), its result is a root
(\texttt{findRootC\_root}) reachable from the start
(\texttt{findRootC\_reachable}), and the checked compression is exactly
\texttt{UF.compress} (\texttt{compressC\_UFArr}), so its partition soundness
transfers. The representation survives the checked compression
(\texttt{UFRep.compress\_preserves}); under a ranked forest
(\texttt{Ranked}) reachable nodes stay in range and the chase stops at a root
within the rank (\texttt{stopsAt\_of\_ranked},
\texttt{findRootC\_ranked\_root}); a bounded acyclic forest gives the same
completion within $n+1$ without a rank (\texttt{stopsAt\_of\_acyclic},
\texttt{findRootC\_acyclic\_root}). A rank implies the acyclicity predicate
(\texttt{acyclic\_of\_ranked}), so proving a rank also discharges the acyclic
route; the concrete 64-byte zero forest is acyclic
(\texttt{uWorld\_acyclic}), and its chase through that route returns 0 with
$n+1$ fuel (\texttt{uWorld\_findRoot\_acyclic}). The construction predicate
\texttt{Built} starts from the identity forest and links a larger-ranked root
under a smaller-ranked one; \texttt{built\_ranked} preserves the rank and
\texttt{forest\_acyclic} derives acyclicity from a construction witness. Orbits are linear, a root reaches only
itself, and compression to a root preserves the partition in both directions
(\texttt{compressC\_same\_iff}); full path compression
(\texttt{UF.compressAll}, \texttt{compressAll\_same\_iff}) preserves the
partition over a supplied path. The path is derived from the chase and
executed through the certified write (\texttt{compressPathCert},
\texttt{compressRootPathCert}); the certificate's final array is
\texttt{compressAll} of its initial array
(\texttt{compressRootPathCert\_compressAll}), so the partition theorem applies
to the executed writes (\texttt{compressRootPathCert\_same\_iff}). Amortized
analysis is not claimed.

\file{Cap/UFCertified.lean} is the certified algorithm-facing interface: a
\texttt{Cert pre} $\alpha$ \texttt{post} is a total computation from a world satisfying
\texttt{pre} to a value and world satisfying \texttt{post}.
\texttt{readParentCert} returns a parent with its range proof, byte bound and
identity as the model's step; \texttt{compressCert} performs the checked write
and preserves \texttt{UFRep} and the invariant; \texttt{findCompressCert} finds
the root and compresses the queried node to it, proving the result a root, the
representation and invariant survive, and \texttt{Same} is preserved in both
directions. \texttt{Cert.pure} and \texttt{Cert.bind} compose computations when
the continuation's precondition follows from the first postcondition, and
\texttt{Cert.bindDep} hands the continuation the first value and its
postcondition with the continuation's post indexed by that value;
\texttt{readThenCompressCert} reads a parent, compresses the same node to the
value read and returns it. The
byte domain is \texttt{n} $\le$ \texttt{256}. \texttt{UFRunner} is the sealed
union-find session (private world, combined invariant, representation, backend and
rank), and \texttt{ufRunner\_findRoot} proves a closed certified execution on the
concrete 64-byte zero forest returns root 0. \texttt{UFRunner.findRootSnapshot}
returns the root with a copy of the allocation's bytes, and
\texttt{ufRunner\_findRootSnapshot} proves the observation is \texttt{(0, 64
zero bytes)}. The compression path is derived from the chase
(\texttt{rootPath}); its nodes are in range and reach the search's result, and
\texttt{rootPath\_compress\_same\_iff} proves full-path partition preservation
along the path the chase takes. The full path is executed through the
certified write: \texttt{compressPathCert} folds \texttt{compressCert} over the
list and \texttt{compressRootPathCert} applies it to \texttt{rootPath}; the
final array is \texttt{compressAll} of the initial array
(\texttt{compressRootPathCert\_compressAll}) and the partition is preserved in
both directions (\texttt{compressRootPathCert\_same\_iff}).

\section{Byte-array backend refinement}

\file{Cap/BackendArr.lean} backs the backend with \texttt{ByteArray} and refines
the pure list model: \texttt{toBackend} reads the array's data back as a byte
list, and empty, reads, writes and range fills agree with \texttt{Backend}
(\texttt{empty\_refines}, \texttt{read\_refines}, \texttt{write\_refines},
\texttt{fills\_refine}). The 64-byte two-range fill through the byte-array route
refines the pure run (\texttt{ba64\_refines}), whose snapshot is the reference
bytes (\texttt{pure64\_snapshot}, kernel-computed). In-place update at runtime
is not observable here and is not claimed; that observation problem is the
subject of the ownership side of this repository.

In one interpreted run on one machine, on the 64-byte fill-and-sum pattern with
agreeing checksums (\texttt{bench/CapBATiming.lean}): 73,959 ns per repetition
for the pure list backend, 47,688 for the byte-array backend, 51,956 for the
reference heap. Interpreted and environment-dependent; no in-place update is
observed or claimed.

\section{Result-returning programs: the context-preserving fragment}

\file{Cap/Typed.lean} gives the context-preserving program interface, repaired:
\texttt{Owned \(\Gamma\) \(\alpha\)} is reads, writes and monadic sequencing over a
fixed permission context, and the read and write constructors take
\texttt{Ctx.CanRead}/\texttt{Ctx.CanWrite} evidence, so an access the context
does not permit is not a program. The evaluator has an explicit outcome
(done/exhausted/refused); \texttt{runOwned\_wf} preserves the world invariant,
\texttt{runOwned\_ctx} leaves the context unchanged, and
\texttt{runOwned\_no\_refusal} proves a typed program is never refused from a
matching, well-formed world. The fragment has no representation index, so a
runtime index cannot acquire evidence there; the data-dependent chase lives in
the certified interface (\file{Cap/UFCertified.lean}), where \texttt{UFRep}
supplies the evidence. It is not a second, weaker access rule. Reads are
observations; allocation, splits, borrow/return and disposal remain in
\file{Cap/Unified.lean}.

\section{The indexed interface: programs with post-contexts}

\file{Cap/IProg.lean} adds the indexed form the handoff asks for next:
\texttt{IProg \(\Gamma\) \(\alpha\) Q} starts with permission context \(\Gamma\),
returns \(a : \\alpha\) and ends with \(Q\\,a\). Context-changing constructors carry their
evidence, so a stale or refused operation is not a program at all: a split that
reuses its parent's capability as a child has an uninhabited freshness proof
(\texttt{stale\_child\_refused}). The fuelled evaluator is total, and the
layer's soundness theorem \texttt{evalIProg\_post} states: execution from a world
matching the typing context ends in a world matching the postcondition,
\texttt{w'.ctx = Q a \(\land\) w'.CommandInv}, so the run keeps well-formedness,
the issuance history and loan readiness. The indexed lifecycle --- allocation, write,
spatial split, write, disposal, with every evidence derived --- is accepted
(\texttt{iLifecycle\_ok}); a second lifecycle adds a fractional split and join
(\texttt{iFracLifecycle\_ok}). It ends on its typed post-context
(\texttt{iLifecycle\_ctx}), well formed (\texttt{iLifecycle\_sound}), with the
pre-existing backend untouched (\texttt{iLifecycle\_snapshot}). Post-contexts
use \texttt{getD} so the typing function stays computable. The loan state is
part of the permission context (\texttt{Ctx.loans}): \texttt{beginBorrow} and
\texttt{endBorrow} carry the owner, leaf, freshness, liveness and loan-readiness
evidence, with success and post-context lemmas; \texttt{iLoanLifecycle} exercises
them and \texttt{iDelegationLifecycle} runs a two-level delegation. The closed
form runs from the empty world: \texttt{iClosed} allocates, writes and disposes,
is accepted, ends with no owners (\texttt{iClosed\_closed}), is well formed, and
leaves the disposed allocation's backend gone (\texttt{iClosed\_freed}).

\section{Historical capability issuance}

The context records `Ctx.issued`, the identifiers it has ever issued. The
world-level spatial split, fractional split and join test a new identifier
against this history as well as against the live set, and extend it on success:
a split that hands a live child the identifier of a retired capability is
refused. \file{Cap/Issuance.lean} carries the regression and its controls: the
prefix of the two-split program is accepted, the reuse program is refused
(\texttt{reuseSplit\_refused}), and the same program with a fresh child is
accepted (\texttt{reuseSplitControl\_ok}); the history is monotone through the
two transitions (\texttt{splitAtChecked\_issued\_sub},
\texttt{joinChecked\_issued\_sub}). Allocation's identifier check and a
history-preservation invariant across every transition remain open, and the
paper states that boundary. The indexed split constructor carries the same
historical-freshness premises the checker tests, and
\texttt{splitAtChecked\_succeeds} proves that a typed split is never refused:
the typed interface and the checker agree on the history check.

\section{Exclusivity of write authority}

Coverage becomes the discipline's central safety property by arithmetic.
\texttt{Tree.budget\_leaf\_ge} shows a present, unconsumed leaf covering a byte
contributes its share to the budget; \texttt{Tree.budget\_two\_leaves} shows two
distinct such leaves contribute both shares. Hence
\texttt{Owner.no\_two\_writable}: at a byte of a well-formed owner, two
writable capabilities have the same identifier --- at most one writer per byte.
The proof needs the leaf's own positivity (the denominator is positive), so the
statement is about well-formed owners, not bare trees.

\section{Access evidence}

\file{Cap/Access.lean} defines the evidence a certified access carries:
\texttt{Ctx.CanWrite} (owner found and the write guard accepts) and
\texttt{Ctx.CanRead} (owner found and the capability readable at the index).
A matching well-formed world supplies the backend and bounds:
\texttt{World.read\_succeeds} and \texttt{World.write\_succeeds} prove
certified accesses are never refused, and \texttt{World.no\_two\_canwrite}
proves two certified writes at one byte name the same capability. In
\file{Cap/IProg.lean} the \texttt{read} and \texttt{write} constructors carry
these predicates, so an unauthorised access is not a program; the lifecycle's
evidence is derived and still decides. Every indexed constructor's evidence now covers its checker:
\texttt{splitAtChecked\_succeeds}, \texttt{World.allocChecked\_succeeds} and
\texttt{World.freeChecked\_succeeds} join the two access success lemmas. The
evaluator's outcome is explicit (\texttt{IOutcome}: completion, fuel exhaustion
or refusal), and \texttt{runIProg\_no\_refusal} lifts the per-command facts to
programs: a typed program is never refused, only completed or exhausted. A
data-dependent index cannot carry write evidence at the program text, so
\file{Cap/UFChecked.lean} supplies the representation predicate that does it at
run time.

\section{The issuance-history invariant}

\file{Cap/History.lean} defines \texttt{Ctx.HistoryOK}: the history is
duplicate-free and every live capability identifier occurs in it. The context
transitions preserve it --- allocation with a fresh, distinct root and disposal
identifier (\texttt{historyOK\_alloc}), disposal (\texttt{historyOK\_free}) and
owner replacement with an unchanged identifier set
(\texttt{historyOK\_update\_of\_ids}, covering a retag) --- with the membership
bounds for both splits and the join. Allocation's checker tests the history, and
the indexed constructor carries the same premises. Each command's history
preservation lives beside its transition: \texttt{historyOK\_write},
\texttt{historyOK\_splitAtChecked}, \texttt{historyOK\_splitFractionChecked},
\texttt{historyOK\_joinChecked}, \texttt{historyOK\_freeChecked},
\texttt{historyOK\_allocChecked}, \texttt{historyOK\_beginBorrow} and
\texttt{historyOK\_endBorrow}. Whole executions preserve it through the command
carrier (\texttt{checkWorld\_top}).

\section{The top-level invariant}

\file{Cap/TopInvariant.lean} combines the world invariant with the issuance
history as \texttt{World.TopInv}. The command carrier
\texttt{World.CommandInv} is exactly \texttt{TopInv} together with
\texttt{Ctx.LoansReady} (\texttt{World.commandInv\_iff}), and every checked
command preserves it: writes, both splits, the join, disposal, allocation, and
borrow/return (\texttt{checkWorld\_top}). Whole executions preserve the carrier,
and the empty and concrete 64-byte worlds satisfy the top-level invariant
(\texttt{emptyWorld\_top}, \texttt{uWorld\_top}). Both split checkers require the split capability to be an active leaf, and the
sibling join refuses non-active children
(\texttt{Tree.joinChildren?\_retired\_none}); the indexed split constructor
carries the active-leaf witness. Splits and fractional joins are refused while
any loan is outstanding on the allocation, so a chain's leaf cannot be
decomposed under a live delegation. The loan state lives in the permission
context (\texttt{Ctx.loans}). Disposal refuses while a loan is attached to the
allocation (\texttt{World.freeChecked\_refused\_of\_loan}), and the indexed
\texttt{free} constructor carries that condition together with the recorded
capability's success. The borrowed-subdivision lifecycle is proved in the
reconstruction model of Section~\ref{sec:reconstruction}; its
integration into this world split/join guard remains open.

\section{The delegation invariant}

\file{Cap/Borrow.lean} states \texttt{World.BorrowOK} as the conjunction of
eleven named predicates: distinct loan, lender, borrower and lifetime names
(\texttt{LoansDistinctOK}); no self-alias (\texttt{LoansNonAliasOK}); issuance of
every name a loan mentions (\texttt{LoansIssuedOK}); live and bounded lifetimes
(\texttt{LifetimesBoundedOK}); an owner at each allocation
(\texttt{LoansAllocOK}); global detachment of every lender from the live
identifier lists (\texttt{LoansDetachedOK}); a recorded chain base and
base-or-chain lender (\texttt{LoansBaseOK}); an active terminal borrower or a
delegated successor with the leaf's resource description
(\texttt{LoansTerminalOK}); adjacent-edge agreement on base, allocation,
footprint, share and ceiling (\texttt{LoansAgreeOK}); issued-order monotonicity
(\texttt{LoansOrderOK}); and the derived \texttt{BorrowAcyclic}. It is preserved
by every command, with \texttt{borrowOK\_beginBorrow} and
\texttt{borrowOK\_endBorrow} for the transitions the mechanism adds. The
world/chain correspondence \texttt{World.delegationChains\_ok} additionally
assumes \texttt{World.LoansChainOrderOK}: \texttt{BorrowOK} does not determine
the order of the world loan list.

\section{Borrowed subdivision and reconstruction}\label{sec:reconstruction}

\file{Cap/Reconstruction.lean} is the component model for the borrowed
subdivision lifecycle; \file{Cap/Integrated.lean} lifts it to a list of
allocations and supplies the sealed boundary. A \texttt{Forest} records at each interior node the
checked spatial or fractional rule that produced its children, and its
\texttt{budget} is the frontier sum (\texttt{budget\_eq\_toTree}); \texttt{Valid}
requires every recorded rule to be checked. \texttt{RuleValid} states the two
rules: a spatial rule needs equal allocation and ceiling, active disjoint
children whose intervals partition the parent, and the inherited share; a
fractional rule needs equal footprint and ceiling, positive child shares, and an
exact sum equal to the parent's. \texttt{Forest.splitSpatialAt} and
\texttt{splitFractionAt} preserve the budget and validity, and
\texttt{reconstructAt} restores a suspended parent from two children matched to
its recorded rule.

\texttt{State} pairs the forest with the reference heap, the issued and retired
names and the outstanding \texttt{Restore} obligations; \texttt{State.Inv} is
frontier coverage, forest validity and the ownership account
\texttt{State.Accounts}. The budget at every byte is still $D$, every recorded
rule is checked, and the account constrains the current authority-node
identities, issuance coverage, retirement, dormant lenders, restoration
realization, resource agreement and restoration order. It is not a structural
restatement of the forest: it relates the forest to \texttt{issued},
\texttt{retired} and \texttt{restores}. \texttt{Restore.Depends} relates two
obligations when the second's lender is the first's borrower, and
\texttt{State.RestoreOrder} requires each dependency to strictly decrease the
index in \texttt{restores}; borrowing prepends, so dependencies run from later
records to earlier ones. \texttt{State.Accounts.dependency\_wellFounded} rules
out a positive delegation cycle, and
\texttt{State.Accounts.restore\_reaches\_terminal} gives every obligation a
finite path to a terminal forest node with resource agreement along the path.
\texttt{Forest.suspended\_frontier\_realizes} connects a suspended parent to its
node entry, rule descendants and descendant frontier budget. A retired or
dormant name is therefore not a current leaf
(\texttt{State.Accounts.retired\_not\_current},
\texttt{State.Accounts.dormant\_not\_current}), and an appended obligation naming
an unissued party breaks the predicate
(\texttt{State.Accounts.cons\_restore\_not\_accounts}). The checked operations are
\texttt{write}, \texttt{borrow}, \texttt{splitSpatial}, \texttt{splitFraction},
\texttt{reconstruct} and \texttt{return}. Each preserves \texttt{State.Inv}
(\texttt{State.write\_inv}, \texttt{State.borrow\_inv},
\texttt{State.splitSpatial\_inv}, \texttt{State.splitFraction\_inv},
\texttt{State.reconstruct\_inv}, \texttt{State.return\_inv}); the split frames
leave the outstanding obligations untouched
(\texttt{splitSpatial\_frames\_restores},
\texttt{splitFraction\_frames\_restores}), a return before reconstruction is
refused (\texttt{return\_none\_of\_no\_leaf}), and a child name that was already
issued or retired is refused (\texttt{splitSpatial\_none\_of\_issued},
\texttt{splitFraction\_none\_of\_retired}).

\file{Cap/Integrated.lean} lifts the state to a list of allocations and supplies
the sealed boundary; \file{Cap/Sealed.lean} re-exports it. The integrated
constructor and invariant proof are private, the checked language is the
\texttt{Cmd} constructors plus a whole-run \texttt{run}, and an accepted call
returns a sealed state with the integrated invariant rebuilt from its
preservation theorem. \texttt{State.Inv} requires distinct live allocation
identifiers, every component's \texttt{State.Inv} (including the ownership
account \texttt{State.Accounts}) and
\texttt{Reconstruction.State.Storage}, and it relates each component's issued
names to a session-persistent capability history. The allocation and capability
histories are append-only: an accepted command only adds names, and
\texttt{free} leaves them intact (\texttt{State.free\_histories\_intact}).
\texttt{State.alloc} refuses an
identifier in \texttt{allocsIssued} and an initial capability name in
\texttt{issued}. A freed identifier therefore cannot be reissued
(\texttt{State.realloc\_after\_free\_refused}) and an old capability cannot
authorize a replacement allocation
(\texttt{State.alloc\_reusing\_old\_capability\_refused}); a fresh identifier with
fresh names is accepted (\texttt{State.alloc\_fresh\_after\_free\_some}).
Observations are a heap
snapshot, byte read, per-fragment authority descriptions, the retired access
names and the outstanding restoration obligations. \texttt{Sealed.lifecycle64}
drives the 64-byte borrow/split/write/reconstruct/return example and
\texttt{Sealed.reuseLifecycle} runs a second full cycle on the first's output
(\texttt{lifecycle64\_ok}, \texttt{lifecycle64\_snapshot}, \texttt{reuse64\_ok}).

The world \texttt{UCmd} path remains separate. \texttt{World.splitAtChecked},
\texttt{splitFractionChecked} and \texttt{joinChecked} still test the
allocation-wide loan exclusion
(\texttt{World.splitAtChecked\_refuses\_with\_loan}), so the world route neither
splits a borrowed range nor an independent sibling range while a loan is
outstanding, and it does not consult this forest. The sealed path executes the
integrated reconstruction language; the completion criterion's single integrated
world-and-restoration execution path is therefore not claimed.

\section{A reusable union-find client}\label{sec:reuse}

\file{Cap/UFReuse.lean} proves the discriminating sequence on the functional
model. From \texttt{[0,0,1,2,2]} the chase from node 3 reads the path
\texttt{[3,2,1]}, and \texttt{reuseFirst\_eq} proves the full-path compression is
exactly \texttt{[0,0,0,0,2]}. The intermediate output is the second operation's
input certificate: \texttt{reuseFirst\_ranked} proves it is still a ranked forest
under the identity rank, and \texttt{reuseFirst\_accepts\_second} supplies the
rank bound, the root and the in-range reachability of every node on the second
path. The chase from node 4 reads \texttt{[4,2]}, and \texttt{reuseSecond\_eq}
proves the second compression is exactly \texttt{[0,0,0,0,0]}. Both preserve the
partition (\texttt{reuseFirst\_same\_iff}, \texttt{reuseSecond\_same\_iff}), and
\texttt{reuseTwoFinds\_same\_iff} composes them back to the starting partition.
\texttt{bench/CapUnionFindReuse.lean} guards the exact arrays. The world-based
sealed \texttt{UFRunner} is still not threaded through the sequence: its world,
representation and rank are private, and it exposes only one find's observation.

\file{Cap/UFIntegrated.lean} runs the two calls through the integrated sealed
boundary on the first call's returned state. A parent array is held in one
allocation's bytes: \texttt{decode} reads a buffer as a \texttt{UF.Arr},
\texttt{Rep s a arr} says the allocation's snapshot decodes to \texttt{arr}, and
\texttt{Rep.fits} bounds every value by a byte. \texttt{chase?} follows parent
pointers with \texttt{Sealed.read?}, \texttt{compressPath} writes the root at
each path node with \texttt{Sealed.write} through the whole-share capability, and
\texttt{sealedFind} returns the root with the successor sealed state.
\texttt{chase?\_eq} proves the sealed chase is the functional
\texttt{UF.findRootF}/\texttt{rootPath} whenever the reads agree with the array.
The first find from node 3 leaves \texttt{[0,0,0,0,2]}; the second, on the
returned state from node 4, leaves \texttt{[0,0,0,0,0]}; a fresh find from node 4
on the start state instead leaves \texttt{[0,0,0,2,0]}
(\texttt{ufSecond\_consumes}). \texttt{ufSecond\_invariant} re-exports the
integrated invariant (including each component's ownership account); a
two-allocation variant is framed as well, with
\texttt{ufSecondOther\_invariant} re-exporting the invariant for the whole state
and \texttt{ufSecondOwnershipFrames\_of\_ok} recording that both allocations'
current fragments are still whole active read--write fragments, no access name
was retired, no return obligation was recorded, both heaps still hold their
bytes, and both allocations are live and distinct. Not proved: a general
per-write refinement or frame law over sealed command lists; the final arrays are
kernel-checked computations of the actual run.

\section{A compiled store composed with an authorized write}\label{sec:kernel}

The separate \file{llvm-bridge/} package closes one kernel-sized instance of the
low-level boundary. Clang compiles
\texttt{void store8(uint8\_t *p, uint64\_t i, uint8\_t v) \{ p[i] = v; \}}; the
emitted body is parsed to an AST that a \texttt{\#guard} pins against the
committed text, and \texttt{Beh\_store8IR\_iff} proves its behaviour at the call
site is exactly the reference store's single bounded outcome under
\texttt{x.toNat < 2\textasciicircum 63}, with no undefined or divergent run.
\file{LlvmBridge/Ownership.lean} then relates one
\texttt{Cap.Integrated} allocation's heap bytes to the bridge's byte memory:
\texttt{allocAddr} maps offset \texttt{off} to \texttt{base + off},
\texttt{MemRep m buf base} says the flat memory holds that buffer at the window,
and \texttt{store8\_exec\_write\_comp} / \texttt{sealed\_store8\_write\_comp}
compose an accepted integrated write with the kernel run: the kernel completes
at every positive budget, the post-memory represents the post-state's updated
buffer, the frame holds outside \texttt{base + off}, the written byte reads back
on both sides, other allocations are unchanged, and the integrated invariant
survives. \texttt{write\_authority} derives the write's current, active,
whole-share, read--write fragment from the command's guard rather than assuming
it.

Stated assumptions, not discharged theorems: the representation premise
\texttt{MemRep m buf base} is supplied by the caller, and nothing assigns
\texttt{base}; \texttt{off < 2\textasciicircum 63} is the pointer-width premise
the emitted \texttt{getelementptr} needs; injectivity is proved within one
window, while disjointness of two allocations' windows is the hypothesis of
\texttt{MemRep.frame\_write}, not a loader theorem; the emitted IR is projected
onto the parser's surface, and parser/AST agreement is a textual guard rather
than a semantic justification of that projection; the capability route is the
integrated reconstruction guard, not the \texttt{Cap.World}
\texttt{Writable}/\texttt{Rel} route; and ABI, native compilation, installation
of the verified body, physical storage and live-handle replay remain trust or
proof obligations (\texttt{Cap.Sealed.boundaryClaims}). In particular this is not
a C memory-safety theorem, a full C/ABI refinement, or a physical ownership
refinement; the bridge's byte memory has no allocation bounds, so the in-bounds
guarantee comes from the ownership guard plus the written offset's range.

\section{Measurements}

\texttt{bench/CapTiming.lean} times the two pure models on the 64-byte
fill/snapshot pattern. At this revision, 10,000 repetitions gave 47,643 ns per
repetition for the backend model and 32,648 ns for the reference model, with equal
checksums. Evaluation is interpreted; this covers the models' operation sequences,
not an interpreter, an optimised backend, FFI, or allocation/GC, and is
environment-dependent.

\section{What is proved, and what is open}\label{sec:open}

\begin{center}
\begin{tabular}{@{}lp{0.64\textwidth}@{}}
\toprule
Proved & Module \\
\midrule
share uniqueness arithmetic, split/join budgets & Share, Footprint, Ledger, Join \\
owner coverage, split/join preservation, guards & Ledger, Context \\
allocation, disposal and identifier freshness & Context \\
loan bookkeeping, delegation invariant, chain correspondence & Loans, Borrow, Discipline \\
nested delegation chain, rooted path, derived holder, checked begin/end & Delegation \\
coupled world/chain state, borrow/return correspondence & Discipline \\
borrowed subdivision, reconstruction and restoration invariant & Reconstruction \\
ownership accounting: issuance, retirement, dormant lenders, restoration realization & Reconstruction \\
integrated sealed execution, persistent histories and reuse refusal & Integrated \\
two-find reusable compression on the functional model & UFReuse \\
two-find compression through the integrated sealed boundary & UFIntegrated \\
compiled store kernel and its composition with one authorized write & llvm-bridge \\
reference heap, snapshot independence & Bytes \\
backend refinement, interpreter agreement, world relation & Backend, Exec, World \\
deep checker soundness, coverage, surface certificate & Core, Surface, Certify \\
union-find and hash-set runs under the guard & UnionFind, HashSet \\
lifetime binding, descendant refusal, borrow/return preservation & Loans, Borrow, Unified \\
one legal-disposal predicate and checked free & Disposal \\
unified command language and world preservation & Unified \\
combined invariant preservation over whole runs & FullInv \\
identifier roles and split freshness & Ledger, Context, Core \\
\bottomrule
\end{tabular}
\end{center}

The two algorithms are owner-layer examples whose guard is not yet load-bearing and
whose store is not tied to the owner by an invariant; they do not run through the
unified execution. That tightening is part of R2 and R1 together.

The principal gap is the connection between certification and execution (review
R1). \texttt{Cap.Unified} now covers the whole command set --- writes, both
splits, join, free, allocation and borrow/return --- with one value-carrying
language, a checker that is also the interpreter, and preservation of the
command carrier, with the lifecycle run through that route. The delegation chain
is repaired as a rooted path with a derived holder, and its checked begin/end
operations preserve validity; the coupled state's borrow and return preserve the
world/chain correspondence. Borrowed subdivision and reconstruction are proved
in the reconstruction model, lifted to the integrated multi-allocation state and
executed through its sealed boundary, with the ownership account and persistent
histories part of the integrated invariant; the two-find
reuse example is proved on the functional model and, for that instance, through
the integrated sealed boundary. One compiled store kernel is proved and composed
with one authorized integrated write at a stated address map. Also open: the read discipline;
storage agreement as a world invariant (R2); integration of the reconstruction
forest into the world split/join guard (R3) --- the reconstruction model exists,
but \texttt{World.splitAtChecked} still refuses every split on a borrowed
allocation and does not consult it, so there is no single integrated
world-and-restoration execution path; threading the world-based sealed
\texttt{UFRunner} through the two-find sequence; one uniform disposal rule as a
predicate with a checked free that is not yet wired into the unified language
(R4); the complete certified 64-byte trace (R5); a general per-write refinement
law over sealed command lists; an address assignment for more than one
allocation, a semantic justification of the IR projection, and the ABI and
installation of the verified kernel; hidden-index elaboration for symbolic
programs; specialisation; and concurrency, which the design defers.

\subsection*{Claim corrections}

\begin{center}
\begin{tabular}{@{}lp{0.52\textwidth}@{}}
\toprule
Earlier wording & Precise statement \\
\midrule
The context is never a public value & The intended runner will encapsulate execution; the current context, session and world are public pure model data \\
Deep checker soundness & Acceptance implies the current guard/transition and coverage results; the unified command language preserves the proved command carrier over whole runs \\
Nested borrowing & The command language carries nested delegation; the chain is a rooted path with a derived holder, and the correspondence and preservation theorems cover chains. Borrowed subdivision is proved in the reconstruction model and executed through the integrated sealed state; the world split/join guard still does not consult it \\
World borrowed reconstruction & The reconstruction model, the integrated multi-allocation state and its sealed boundary prove subdivision, reconstruction and restoration; \texttt{World.splitAtChecked} keeps the allocation-wide loan exclusion and does not consult the forest, so no single integrated world/sealed path is claimed \\
64-byte acceptance run & The unified lifecycle runs alloc, write, split, write, fractional split, join, write and free; the integrated sealed lifecycle runs borrow, split, write, reconstruct and return in the checked multi-allocation state. A single world trace combining a live borrow with decomposition awaits integration of the reconstruction forest into the world guard \\
Mutable backend & A pure buffer state machine over lists, not native mutation or RC=1 \\
Admission excludes closures & Canonical storage excludes arbitrary original values; the encoding relation alone is not confinement \\
Reconstruction invariant & Frontier coverage and rule validity alone are not the account; \texttt{State.Inv} also carries \texttt{State.Accounts}, which relates current identities to the issued, retired and outstanding restoration histories and rejects malformed obligations \\
Allocation reuse & Prevention is a session-persistent allocation and capability history, not a generation or a physical-slot check; replaying an old sealed state is outside the model \\
Compiled store & One clang-emitted kernel is proved and composed with one authorized integrated write at a caller-supplied address window. This is not a C memory-safety theorem, a full C/ABI refinement, a loader address assignment, a semantic justification of the IR projection, or a physical ownership refinement \\
\bottomrule
\end{tabular}
\end{center}

\section{How to run}

\begin{lstlisting}[style=out]
nix develop -c lake build Cap
nix develop -c bash scripts/check_all.sh

# examples
nix develop -c lake env lean bench/NatThunkAdmission.lean
nix develop -c lake env lean bench/CapNegative.lean
nix develop -c lake env lean bench/CapSurfaceReject.lean   # intended refusals
nix develop -c lake env lean bench/CapTiming.lean

# positive and negative capability fixtures, with their intended diagnostics
nix develop -c bash scripts/check_cap_controls.sh

# borrowing, reconstruction and reuse axiom audit
nix develop -c lake env lean scripts/check_cap_axioms.lean
# live-documentation claim census, read from the working tree
bash scripts/check_cap_docs.sh --rev WORKTREE
\end{lstlisting}

The harness is part of \texttt{check\_all}: the positive fixtures must elaborate,
and each negative fixture must fail with its intended diagnostic rather than an
unrelated error.

The \file{llvm-bridge/} composition is a separate package with a sibling
\file{lean-llvm} dependency; \file{llvm-bridge/scripts/build-remote.sh} builds it
on a worker and \file{llvm-bridge/EVIDENCE.md} records those builds. No local
Lean build of it was run for this revision.

The LaTeX outputs were not rebuilt at this revision because the worker image has
no LaTeX installation: \file{report/capability-report.pdf},
\file{report/cap-tutorial.pdf} and \file{paper/sealed-ownership.pdf} are older
than their \file{.tex} sources. The sources are updated; the PDFs are not.

\end{document}
